% ============================================================================
% Fichier cible : ARIEL_LOGMINER_MEMOIRE_FINAL/chapters/chapitre1_introduction.tex
%   (Introduction générale, « Positionnement et définition opérationnelle des
%   termes », PDF p. 4)
% Passage à remplacer : la première phrase du premier paragraphe de la section,
%   intégrée au lot A1-a : « Le titre du mémoire qualifie l'architecture de
%   \textit{multi-agents} et de \textit{légère} ; la question centrale y ajoute
%   les qualificatifs \textit{autonome}, \textit{distribué} et
%   \textit{multitâche}, et l'implémentation désigne ses agents comme
%   \textit{intelligents}. » La phrase suivante (« Ces termes peuvent couvrir
%   des réalités très différentes [...] ») est conservée.
% Auteur : Agent 2, chantier 2026-11-05.
% Motif : à n'intégrer QUE si la question centrale reformulée
%   (chapitre1_introduction.tex__question_centrale_QR.tex) est retenue ; elle
%   ne contient plus « autonome », « distribué » ni « multitâche ». Les
%   définitions de la section restent nécessaires : les agents sont qualifiés
%   d'autonomes, d'intelligents et de multitâches dans le corps (§2, §3, §6.3 ;
%   classe MultiTaskIntelligentAgent, chapitre4_implementation.tex).
% ============================================================================

Le titre du mémoire qualifie l'architecture de \textit{multi-agents} et de \textit{légère} ; le mémoire décrit en outre ses agents comme \textit{autonomes}, \textit{intelligents} et \textit{multitâches}, et leur exécution comme \textit{distribuée}.
